% ID: FIS-OTT-RIF-003
% Titolo: Spostamento laterale in una lastra di vetro
% Disciplina: Fisica
% Area: Ottica
% Argomento: Rifrazione
% Sottoargomento: Lastra a facce parallele
% Classe: Seconda liceo scientifico
% Difficolta: 4
% Tipo: problema_numerico
% Risultato: soluzione da aggiungere
% Tempo_stimato: 12 min
% Competenze: applicazione della legge di Snell, geometria del raggio, modellizzazione
% Prerequisiti: rifrazione, legge di Snell, trigonometria
% Tag: fisica, ottica, rifrazione, lastra_vetro, spostamento_laterale
% Fonte: verifica_fisica_ottica_anonimizzata
% Autore: Riccardo Carli
% Licenza: CC-BY-SA-4.0

\begin{esercizio}
Un raggio di luce attraversa una lastra di vetro piana a facce parallele,
immersa in aria. L'angolo di incidenza sulla prima superficie e
\(30{,}0^\circ\), lo spessore della lastra e \(6{,}00\,\mathrm{mm}\) e
l'indice di rifrazione del vetro e \(n=1{,}52\).

Calcola lo spostamento laterale del raggio emergente rispetto alla direzione
del raggio incidente.
\end{esercizio}

\begin{soluzione}
% Soluzione da aggiungere.
\end{soluzione}

